\documentclass[../main.tex]{subfiles}
\begin{document}

Ransomware has become one of the most serious threats to organizations of
all types. Modern attacks do not merely encrypt the victim's data: they
systematically destroy backup copies before encryption, eliminating the
primary recovery path. In response, the industry proposes immutable storage
as a protection mechanism, yet the academic literature lacks experimental
comparisons that quantify its actual impact on recovery time.

This Bachelor's Thesis presents the design, implementation and evaluation
of a ransomware simulation laboratory based entirely on Docker containers.
The environment replicates a minimal enterprise infrastructure consisting
of a file server, a workstation, an SFTP backup server, an object storage
system with Object Lock protection (MinIO), and a complete SIEM stack
(Wazuh~4.7.5 with OpenSearch).

The experimental design follows a factorial scheme of four attack scenarios,
three recovery strategies and ten repetitions, totalling 120~fully automated
runs. The scenarios cover the most relevant tactics of the MITRE ATT\&CK
matrix: mass file encryption (T1486), active destruction of backup copies
(T1490), lateral movement via SSH (T1021.004), and data exfiltration
(T1041). The three strategies evaluated are: no backup, a Restic repository
over conventional SFTP, and a Restic repository over MinIO with COMPLIANCE
Object Lock active for 30~days.

Results show that the primary hypothesis of at least a 30\,\% MTTR reduction
is not statistically confirmed: when the repository is available, both
strategies restore at equivalent speed. However, the active destruction
scenario reveals the fundamental difference: the SFTP strategy fails in
100\,\% of cases as the repository is destroyed, while Object Lock blocks
all deletion attempts and maintains a 100\,\% success rate (Fisher exact
test $p < 0.001$). The advantage of immutable storage is not speed but
guaranteed availability.

Detection times via Wazuh range from 5.7 to 41.9~seconds, compared to the
more than 200~days reported for manual detection, confirming that the SIEM
qualitatively transforms incident response capability.

As an additional contribution, the laboratory is validated using the Atomic
Red Team framework, executing all four ATT\&CK techniques against the live
environment. All are correctly detected by the developed rules. The
laboratory is published as free software to facilitate reproduction.

\end{document}
